% Options for packages loaded elsewhere
\PassOptionsToPackage{unicode}{hyperref}
\PassOptionsToPackage{hyphens}{url}
\PassOptionsToPackage{dvipsnames,svgnames,x11names}{xcolor}
\documentclass[
]{ltjsarticle}
\usepackage{xcolor}
\usepackage[margin=22mm]{geometry}
\usepackage{amsmath,amssymb}
\setcounter{secnumdepth}{-\maxdimen} % remove section numbering
\usepackage{iftex}
\ifPDFTeX
  \usepackage[T1]{fontenc}
  \usepackage[utf8]{inputenc}
  \usepackage{textcomp} % provide euro and other symbols
\else % if luatex or xetex
  \usepackage{unicode-math} % this also loads fontspec
  \defaultfontfeatures{Scale=MatchLowercase}
  \defaultfontfeatures[\rmfamily]{Ligatures=TeX,Scale=1}
\fi
\usepackage{lmodern}
\ifPDFTeX\else
  % xetex/luatex font selection
\fi
% Use upquote if available, for straight quotes in verbatim environments
\IfFileExists{upquote.sty}{\usepackage{upquote}}{}
\IfFileExists{microtype.sty}{% use microtype if available
  \usepackage[]{microtype}
  \UseMicrotypeSet[protrusion]{basicmath} % disable protrusion for tt fonts
}{}
\makeatletter
\@ifundefined{KOMAClassName}{% if non-KOMA class
  \IfFileExists{parskip.sty}{%
    \usepackage{parskip}
  }{% else
    \setlength{\parindent}{0pt}
    \setlength{\parskip}{6pt plus 2pt minus 1pt}}
}{% if KOMA class
  \KOMAoptions{parskip=half}}
\makeatother
\usepackage{longtable,booktabs,array}
\newcounter{none} % for unnumbered tables
\usepackage{calc} % for calculating minipage widths
% Correct order of tables after \paragraph or \subparagraph
\usepackage{etoolbox}
\makeatletter
\patchcmd\longtable{\par}{\if@noskipsec\mbox{}\fi\par}{}{}
\makeatother
% Allow footnotes in longtable head/foot
\IfFileExists{footnotehyper.sty}{\usepackage{footnotehyper}}{\usepackage{footnote}}
\makesavenoteenv{longtable}
\setlength{\emergencystretch}{3em} % prevent overfull lines
\providecommand{\tightlist}{%
  \setlength{\itemsep}{0pt}\setlength{\parskip}{0pt}}
\usepackage{bookmark}
\IfFileExists{xurl.sty}{\usepackage{xurl}}{} % add URL line breaks if available
\urlstyle{same}
\hypersetup{
  colorlinks=true,
  linkcolor={Maroon},
  filecolor={Maroon},
  citecolor={Blue},
  urlcolor={Blue},
  pdfcreator={LaTeX via pandoc}}

\author{}
\date{}

\begin{document}

\section{総説:関係ひとつ・定数ひとつ・1ビットから ---
標準理論への辞書と公理の交換収支}\label{ux7dcfux8aacux95a2ux4fc2ux3072ux3068ux3064ux5b9aux6570ux3072ux3068ux30641ux30d3ux30c3ux30c8ux304bux3089-ux6a19ux6e96ux7406ux8ad6ux3078ux306eux8f9eux66f8ux3068ux516cux7406ux306eux4ea4ux63dbux53ceux652f}

\textbf{著者}: 木原範昭 \textbf{作成}: 2026年6月12日 \textbf{版}:
v1.0(確定版・公開対応。論文13〜16 の Zenodo 公開に伴い出典 DOI
を確定、§5 状態列を全「公」化、引用地図の機械照合完了{[}Claude
Code{]}。スナップショット DOI
取得・記入済み{[}10.5281/zenodo.20666114{]}。claude.ai
起草、DOI手続き・tex/PDF=Claude Code) \textbf{DOI(本版)}:
10.5281/zenodo.20666133 / \textbf{Concept DOI}: 10.5281/zenodo.20666132
\textbf{ライセンス}: CC BY 4.0 \textbf{シリーズ}:
波長空間と周波数空間の双対幾何(総説。番号未定:論文17候補/独立総説/note
原型)

\begin{center}\rule{0.5\linewidth}{0.5pt}\end{center}

\subsection{要旨}\label{ux8981ux65e8}

本シリーズの出発点は、関係式ひとつ(νλ=1)、定数ひとつ(零点
½)、そして文字通りの1ビット(非対称1ビット)である ---
公理の情報量はこれで全部であり、これに状態の読み方を定める二つの運用原理(配置読み・記録原理)が付く。\textbf{持ち込まなかったもの}:ゲージ原理、複素数体とその共役、時空次元、時間パラメータ、対称化公準、測定公準、Born
則、作用。ここまでのシリーズは、誤読と否認を避けるために\textbf{物理的な接続をあえて語らず}、段階的に検算可能な体系として論文を積み重ねてきた。しかし当初からの目標は一貫して、この極めて限られた仮定だけから、標準量子理論・標準ゲージ理論に近い豊かな表現力を持つ系を導出できるかを確かめることであった。本総説はその目標に対する現況の収支報告である。方法は再導出ではなく、\textbf{項目検証された対応辞書}と\textbf{移送原理}。主結果:(1)
運動学的辞書は埋まり、全項が機械検証済み・\textbf{全出典が公開済み}である。(2)
\textbf{標準理論が公理として置く五項 ---
対称化・時間・測定・次元・電荷保存 --- を、本体系は定理として持つ}。(3)
未記入項は二つ ---
測度と作用。さらに本稿は\textbf{不完全接続点そのものを解剖し(§9)、その構造的原因を特定する}。完全に埋まった辞書は再記法を意味し、経験的内容を持たない
---
\textbf{接続が不完全な点こそ、両理論が異なる主張をし、観測が裁ける唯一の場所}であり、残額は欠陥であると同時に本体系の判定可能性の座である。\textbf{結論}:当初は単なるトイモデルを目指していたが、辞書(§5)を確認いただきたい
---
標準量子理論やゲージ理論の基本要件に、かなり近い水準で接続できたと考える。\textbf{注意}:本シリーズは新たな理論物理学や統一理論を目指す試みでは決してない。\textbf{別のシンプルな仮定からでも、標準量子理論に近い構造が導出できるかもしれない}
---
そういうモデルとして読んでいただきたい。物理的同一視は行わない。追記(§13)として、出発点の経緯の総括と、「これは単なる複素数の言い換えではないか」という予想される批判への回答を収載する。全出典への参照とアクセス方法は参考文献節に与える
--- 本稿の主張は内容だけでなく\textbf{過程においても}監査可能である(§2
末尾)。

\begin{center}\rule{0.5\linewidth}{0.5pt}\end{center}

\subsection{1.
読み方の注意(最初に)}\label{ux8aadux307fux65b9ux306eux6ce8ux610fux6700ux521dux306b}

本総説とシリーズ全体は、次のものでは\textbf{ない}:新たな理論物理学の提案、統一理論の試み、標準理論への対抗仮説。標準理論は厳密に正しく、その地位を本シリーズは一切争わない。

本シリーズが確かめようとしたのは一つの問いである ---
\textbf{標準量子理論の構造は、どれほど少ない・どれほど異質な仮定から再出現しうるか。}
これは量子論の再構成プログラムと同じ形の問いだが、本シリーズの特徴は仮定の情報量が極端に小さいこと(§3:関係1・定数1・1ビット)と、全段が機械検算可能な離散系で進むことにある。したがって本稿の主張が正しくても標準理論は何も失わず、間違っていれば本稿だけが失う。その意味で安全に読める文書である。「別のシンプルな仮定からでも標準量子理論に近い構造が導出できるかもしれないモデル」---
これが著者の指定する読み方であり、シリーズ全篇の免責「物理的同一視は行わない」はこの読み方の技術的表現である。

\subsection{2. 経緯 ---
なぜ今まで接続を語らなかったか}\label{ux7d4cux7def-ux306aux305cux4ecaux307eux3067ux63a5ux7d9aux3092ux8a9eux3089ux306aux304bux3063ux305fux304b}

ここまでの論文1〜16と補遺84本は、物理的な接続の主張を意図的に避けてきた。理由は二つある。\textbf{誤読の回避}:接続を先に語れば、未検証の同一視と混同される。\textbf{否認の回避}:検証不能な主張は門前払いされる。そこで戦略を二期に分けた
---
第一期は、接続を語らずに、段階的・機械検算可能・二者独立検算つきの体系だけを積む(完了)。第二期が本総説である:積み上がった辞書を根拠に、接続の現況を初めて正面から勘定する。当初の目標(限られた仮定からの豊かな表現力の導出)は第一期の全期間を通じて不変だった
--- 語らなかっただけである。

\textbf{記録の所在(過程の監査可能性)}:本総説の主張は内容だけでなく過程においても監査できる。第一期の全過程
---
補遺84本、全検証スクリプト、二者独立検算の往復、そして\textbf{予言の反証・説明の撤回・定義の照合を含む相互訂正の履歴}、さらに論文13〜16
の査読記録(各2ラウンド・計8査読)---
は公開リポジトリ(参考文献)にコミット履歴つきで残されている。すなわち本稿の表は突然の思いつきではなく、\textbf{日付のある思考実験と機械検証の累積}の収支であり、その累積自体が第三者の検証対象として公開されている。失敗と訂正を含む履歴を消さずに残すこと
--- それが本シリーズの検算文化であり、本稿の信頼性の根拠である。

\subsection{3. 仮定の勘定 ---
スタートの軽さこそ第一の結果}\label{ux4eeeux5b9aux306eux52d8ux5b9a-ux30b9ux30bfux30fcux30c8ux306eux8efdux3055ux3053ux305dux7b2cux4e00ux306eux7d50ux679c}

\subsubsection{3.1
公理(三つ、情報量で数えれば約2.5)}\label{ux516cux7406ux4e09ux3064ux60c5ux5831ux91cfux3067ux6570ux3048ux308cux3070ux7d042.5}

{\def\LTcaptype{none} % do not increment counter
\begin{longtable}[]{@{}lll@{}}
\toprule\noalign{}
\# & 公理 & 情報量 \\
\midrule\noalign{}
\endhead
\bottomrule\noalign{}
\endlastfoot
1 & 相反双対 νλ=1 & 関係式ひとつ \\
2 & 零点 ½ & 定数ひとつ \\
3 & 非対称1ビット & \textbf{文字通り1ビット} \\
\end{longtable}
}

第三公理は二値の選択であり、公理として可能な最小の入力である。「仮定は実質
2.5 件」という勘定はこの情報量の正確な表現である。

\subsubsection{3.2
運用原理(二つ)}\label{ux904bux7528ux539fux7406ux4e8cux3064}

\textbf{配置読み}(状態=占有セルの集合=波そのもの)と\textbf{記録原理}(記録は無→有の追記、存在する記録は現在読める)。これらは内容の公理ではなく読み方の規則だが、批判的な勘定では数えるべきであり、数える。重要なのは、これらが標準理論側の\textbf{状態公準・測定公準の対応物を兼ねる}こと
--- 追加コストではなく交換の一部である(§6)。

\subsubsection{3.3
持ち込まなかったもののリスト}\label{ux6301ux3061ux8fbcux307eux306aux304bux3063ux305fux3082ux306eux306eux30eaux30b9ux30c8}

ゲージ原理。複素数体 ℂ
とその共役操作。時空の次元。外部時間パラメータ。対称化公準。測定公準(射影・状態更新)。Born
則。作用・ラグランジアン。---
いずれも公理に含まれない。このうちゲージ骨格・複素構造・Z₄
位相・次元・時間・対称化・測定は\textbf{体系内で出現する}(§5, §7)。Born
と作用が残額である(§8, §9)。

\subsection{4.
方法:辞書と移送原理}\label{ux65b9ux6cd5ux8f9eux66f8ux3068ux79fbux9001ux539fux7406}

\textbf{辞書}は対象の対 \{体系の対象 ↔ 標準理論の対象\}
の表であり、各項は (i) 体系側の構成が定理または機械検証として確立、(ii)
構造的対応(代数・順序・変換則)の確認、をもって「埋まった」とする。\textbf{移送原理}:標準理論の定理
T の前提が全て辞書の埋まった項に対応するとき、T
は再導出なしに移送される。再導出を選ばない理由は根本的である:標準理論の証明は正しい。必要なのは証明のやり直しではなく、前提の対応物の検証である。座標写像(論文15)の完成により、移送は「原理」から「演算」になった。辞書は物理的同一視を主張しない
---
主張するのは構造対応であり、同じ結果を出す別の写像は、同一視なしに移送を支える。

\subsection{5. 辞書の現況(主表 ---
結論の根拠はこの表である)}\label{ux8f9eux66f8ux306eux73feux6cc1ux4e3bux8868-ux7d50ux8ad6ux306eux6839ux62e0ux306fux3053ux306eux8868ux3067ux3042ux308b}

凡例 --- 状態:\textbf{公}=Zenodo
公開済み(\textbf{本稿時点で全出典が公開済み} --- DOI
は参考文献)。検証:定理(証明あり)/機検(機械検証)。

{\def\LTcaptype{none} % do not increment counter
\begin{longtable}[]{@{}
  >{\raggedright\arraybackslash}p{(\linewidth - 6\tabcolsep) * \real{0.2500}}
  >{\raggedright\arraybackslash}p{(\linewidth - 6\tabcolsep) * \real{0.2500}}
  >{\raggedright\arraybackslash}p{(\linewidth - 6\tabcolsep) * \real{0.2500}}
  >{\raggedright\arraybackslash}p{(\linewidth - 6\tabcolsep) * \real{0.2500}}@{}}
\toprule\noalign{}
\begin{minipage}[b]{\linewidth}\raggedright
体系の対象
\end{minipage} & \begin{minipage}[b]{\linewidth}\raggedright
標準理論側の対応
\end{minipage} & \begin{minipage}[b]{\linewidth}\raggedright
状態
\end{minipage} & \begin{minipage}[b]{\linewidth}\raggedright
出典
\end{minipage} \\
\midrule\noalign{}
\endhead
\bottomrule\noalign{}
\endlastfoot
占有セル集合(配置読み) & 状態 & 公・定理+機検 & 論文7 \\
波の辞書・R 量子化 & 状態↔波動の対応 & 公・機検 & 論文5 \\
Σν² 保存 & エネルギー様保存量 & 公・定理+機検 & 論文6 \\
排他(二重占有禁止) & 排他原理の対応物 & 公・定理 & 論文7 \\
内部時間 t=Σ1/√s & 時間パラメータ & 公・\textbf{定理(導出量)} &
論文8、論文13 \\
位置=符号セクター(状態追加ゼロ) & 位置観測量 & 公・定理+機検 & 論文13 \\
空間=保存量格子の Pontryagin 双対 & 位置表示 & 公・定理 & 論文13 \\
複素構造 C⊂H(時間方向の選択から) & 複素ヒルベルト空間の i & 公・定理 &
論文14 \\
導出 Z₄ 輸送位相 & 位相 / U(1) の骨格 & 公・定理+機検 & 論文14 \\
正準規約定理(同種対称化) & \textbf{対称化公準} & 公・定理 & 論文14 \\
次元4の必然性(\{1,2,4,8\}×平方数) & \textbf{時空次元(入力)} & 公・定理 &
論文11 \\
測定=追記+記録定理 & \textbf{測定公準} & 公・定理+機検 & 論文12 \\
電荷様パリティ選抜則 & \textbf{電荷保存} & 公・定理+機検 & 論文15
定理8 \\
座標写像 xyztRQ(順・逆・保存可換) & 座標系 & 公・機検(全点独立再現) &
論文15 \\
ゲージ骨格(計数から創発) & ゲージ構造(\textbf{作用は未構成}) & 公・機検
& 論文10 \\
等速運動の運動学 & 自由運動の運動学 & 公・機検 & 論文15 \\
\end{longtable}
}

\subsection{6. 公理の交換収支 ---
標準理論が公理で置く五項を、本体系は定理として持つ}\label{ux516cux7406ux306eux4ea4ux63dbux53ceux652f-ux6a19ux6e96ux7406ux8ad6ux304cux516cux7406ux3067ux7f6eux304fux4e94ux9805ux3092ux672cux4f53ux7cfbux306fux5b9aux7406ux3068ux3057ux3066ux6301ux3064}

{\def\LTcaptype{none} % do not increment counter
\begin{longtable}[]{@{}
  >{\raggedright\arraybackslash}p{(\linewidth - 4\tabcolsep) * \real{0.3333}}
  >{\raggedright\arraybackslash}p{(\linewidth - 4\tabcolsep) * \real{0.3333}}
  >{\raggedright\arraybackslash}p{(\linewidth - 4\tabcolsep) * \real{0.3333}}@{}}
\toprule\noalign{}
\begin{minipage}[b]{\linewidth}\raggedright
標準理論側の地位
\end{minipage} & \begin{minipage}[b]{\linewidth}\raggedright
本体系での地位
\end{minipage} & \begin{minipage}[b]{\linewidth}\raggedright
出典
\end{minipage} \\
\midrule\noalign{}
\endhead
\bottomrule\noalign{}
\endlastfoot
対称化公準(同種粒子) & \textbf{定理}(正準規約定理) & 論文14 \\
時間=外部パラメータ & \textbf{定理}(t
は導出量:独立帳簿なし・独立逆写像なし) & 論文13・15 \\
測定公準(射影・更新) & \textbf{定理}(追記+記録定理) & 論文12 \\
時空次元=入力 & \textbf{定理}(必然性) & 論文11 \\
電荷保存(対称性を入力として) & \textbf{定理}(台のパリティ算術 ---
対称性の入力なし) & 論文15 定理8 \\
\end{longtable}
}

支払い:関係1+定数1+1ビット+運用原理2。受取:上記五項に加え、複素構造・Z₄
位相・ゲージ骨格・座標写像。Born
則は別勘定(§8.1):輸入という退路(輸入しても公理勘定で標準理論と同水準に並ぶ)と、定理化という上振れの両方が開いている。

\subsection{7.
「持ち込まなかったものが出てくる」各個撃破の記録}\label{ux6301ux3061ux8fbcux307eux306aux304bux3063ux305fux3082ux306eux304cux51faux3066ux304fux308bux5404ux500bux6483ux7834ux306eux8a18ux9332}

\begin{enumerate}
\def\labelenumi{\arabic{enumi}.}
\tightlist
\item
  \textbf{位置が読めない} →
  位置=符号セクター。状態の追加ゼロで解決(論文13)。
\item
  \textbf{複素数がない} → 時間方向の選択が複素構造 C⊂H
  を選ぶ。複素共役も公理でなく、反転対合として内部に出現(論文14)。
\item
  \textbf{位相の出所} → Z₄
  輸送位相は導出される(論文14)。向き付けゲージ固定の必然性も同時に判明。
\item
  \textbf{同種粒子の扱い} →
  正準規約定理(規約の恣意性が定理に変わる、論文14)。
\item
  \textbf{移送原理自体の障害}(監査指摘) → 条件の明示と昇格で解消(論文14
  に収載)。
\item
  \textbf{座標への着地} → 論文15
  が5段の演算として構成・機械検証(往復500/500、保存可換)。
\end{enumerate}

\subsection{8.
未記入の二項(辞書の空欄)}\label{ux672aux8a18ux5165ux306eux4e8cux9805ux8f9eux66f8ux306eux7a7aux6b04}

\subsubsection{8.1 測度}\label{ux6e2cux5ea6}

体系は経路振幅を自由パラメータなしで与える ---
確率の素材は揃っている。決まっていないのは集計規則で、有限の決定構造に分解済み:\textbf{干渉粒度
× イベント構造1ビット ×
枝チャネル条件付け}(論文16:選定問題の厳密な縮約と判定表 ---
公開済み、判定実験は進行中)。下限(退路):Born
則を公理として輸入すれば公理勘定で標準理論と並ぶ ---
接続の下限は確保済み。上振れ:記録原理が集計を一意化すれば Born
は定理になる。合格基準は三つ組 --- 収束(然るべき極限で Born
に合流)・抑制(検証済み領域で偏差が統計的検閲 \textasciitilde1/√N
以下)・予言(深い離散領域での検証可能な偏差)。

\subsubsection{8.2 作用・力学}\label{ux4f5cux7528ux529bux5b66}

未構成である(論文10
の自己申告のまま)。等速運動の運動学(論文15)はあるが、「力なし→等速」の導出・相互作用の力学はない。

\subsection{9. 不完全接続の解剖 ---
残額はなぜ残るか、そして残ることの何が鍵か}\label{ux4e0dux5b8cux5168ux63a5ux7d9aux306eux89e3ux5256-ux6b8bux984dux306fux306aux305cux6b8bux308bux304bux305dux3057ux3066ux6b8bux308bux3053ux3068ux306eux4f55ux304cux9375ux304b}

接続の不完全な箇所を隠さず、\textbf{解剖する}。まず分類する。

\subsubsection{9.1 二つの種類}\label{ux4e8cux3064ux306eux7a2eux985e}

\begin{itemize}
\tightlist
\item
  \textbf{工事中}:未完だが、完成が原理的に可能な箇所(同型定理、連続極限写像の構成)。これは時間の問題である。
\item
  \textbf{構造的不一致}:完全一致が\textbf{原理的に期待されない}箇所。本節の主題はこちらである。
\end{itemize}

\subsubsection{9.2 解剖表}\label{ux89e3ux5256ux8868}

{\def\LTcaptype{none} % do not increment counter
\begin{longtable}[]{@{}
  >{\raggedright\arraybackslash}p{(\linewidth - 6\tabcolsep) * \real{0.2500}}
  >{\raggedright\arraybackslash}p{(\linewidth - 6\tabcolsep) * \real{0.2500}}
  >{\raggedright\arraybackslash}p{(\linewidth - 6\tabcolsep) * \real{0.2500}}
  >{\raggedright\arraybackslash}p{(\linewidth - 6\tabcolsep) * \real{0.2500}}@{}}
\toprule\noalign{}
\begin{minipage}[b]{\linewidth}\raggedright
箇所
\end{minipage} & \begin{minipage}[b]{\linewidth}\raggedright
不完全性の内容
\end{minipage} & \begin{minipage}[b]{\linewidth}\raggedright
構造的原因
\end{minipage} & \begin{minipage}[b]{\linewidth}\raggedright
反転価値(判定鍵としての地位)
\end{minipage} \\
\midrule\noalign{}
\endhead
\bottomrule\noalign{}
\endlastfoot
\textbf{測度(Born)} &
集計規則が未選定。候補測度(チャネル一貫/配置粒度)は有限 s
で\textbf{厳密に相違}(鏡像親で W=0 ∧ W′\textgreater0 等 ---
数値誤差でなく代数的構造) & 標準理論側にも選定機構はない(Born
は公理)。本体系は Born
を\textbf{写し取ることを拒否する構造}を内蔵している:複数候補が有限離散で真に分岐し、合流できるのは高々一系
& 本体系の経験的内容の最大の在処。Born
が厳密法則なら上振れは閉じる(退路で接続維持)。Born
が\textbf{極限有効則}なら、深い離散領域に本体系型(粒度・パリティ・鏡像相殺構造)の偏差が住む
--- 偏差の「形」まで指定できるのは選定機構を持つ側だけである \\
\textbf{作用・力学} & 未構成 &
\textbf{移送原理の前提検査で止まる}:作用は外部時間 t
上の積分として定式化されるが、辞書の t の欄には「導出量」(独立の軸なし
---
定理)が入っている。前提の対応物が\textbf{別の地位}を持つため、そのままの移送は原理的に不可。力学は記録列の言語への再定式化を経てのみ接続できる
& 標準理論自身が重力界面で\textbf{同じ困難}(時間の問題:外部 t
の不在)に遭うことが知られている。本体系が強いられている再定式化 ---
力学=記録の簿記 ---
は、標準理論があの界面で強いられるはずの再定式化と同型でありうる。接続の失敗点が、標準理論自身の拡張失敗点と一致している \\
\textbf{離散 ↔ 連続(横断)} &
辞書の全行に通底:本体系は有限・離散、標準理論は連続・無限次元。接続は対応と包含までで、極限写像は未構成
&
離散性は近似の選択でなく\textbf{公理}(格子・1ビット)。連続側に合流するか否かは構成課題であり、前提でない
&
「どちらが基本でどちらが近似か」という判定問題そのものを立てる。深部に離散性の痕跡(統計的検閲
\textasciitilde1/√N、粒度構造)が見つかれば連続記述が有効側。見つからなければ本体系が有効模型側。\textbf{判定は双方向に可能} \\
\textbf{連続対称性} & 体系の対称性は B₄(離散)と Z₄(導出位相)。U(1)
全体・連続回転との接続は包含(Z₄↪U(1))まで &
離散公理の直接の帰結(横断不一致の派生) & 位相の Z₄
量子化・深部異方性は、連続対称性が厳密なら存在しない構造。識別可能な構造的予言の座 \\
\textbf{結合の力学} & ゲージ骨格はあるが結合定数の力学はない &
作用未構成と同根。加えて標準理論側では結合定数(例:α)は\textbf{測定入力
--- 対応する定理の欄が標準側に存在しない} &
接続の不在がそのまま「\textbf{本体系側にしか欄が書けない}」形になっている唯一の例。計数側には恒等式候補が存在する(別系列・Zenodo
公開済みの α⁻¹=137+(π²/2)α 型恒等式 ---
シリーズ外参照として注記)。標準理論が原理的に計算できない量に、計数体系は候補式を持つ \\
\end{longtable}
}

\subsubsection{9.3
反転価値の論理(本節のテーゼ)}\label{ux53cdux8ee2ux4fa1ux5024ux306eux8ad6ux7406ux672cux7bc0ux306eux30c6ux30fcux30bc}

仮に辞書が完全に埋まり尽くしたとせよ。そのとき本体系は標準理論の\textbf{再記法}であり、いかなる実験も両者を区別できない
---
学術的価値(公理の最小性の実証)は残るが、経験的内容は消える。\textbf{したがって、本体系が標準理論に対して何かを「言える」場所は、接続が不完全な点に限られる。}
残額(測度・作用)と横断不一致(離散/連続)は欠陥の表であると同時に、両理論が異なる主張をする点の完全な一覧
---
すなわち\textbf{観測が将来どちらかに軍配を上げうる場所の地図}である。

判定の方向は双方向に明示できる。離散性の痕跡・Born
からの構造化された偏差・記録的力学の必要が観測される未来では、本体系の「不完全な接続」は不完全だったのではなく、\textbf{標準理論の側が極限有効則だった}ことの先行記述になる。逆に、Born
厳密・連続基本・外部時間で全てが押し切れる未来では、本体系は有効模型ないし再記法に格下げされる。どちらに転んでも判定が立つこと
--- それが、この不完全性を隠さず解剖して公開する理由である。

\subsubsection{9.4
規律(本節が主張しないこと)}\label{ux898fux5f8bux672cux7bc0ux304cux4e3bux5f35ux3057ux306aux3044ux3053ux3068}

本節は偏差の\textbf{実在}を主張しない。主張するのは偏差が住むべき\textbf{構造の所在}(粒度・パリティ・量子化・記録性)の特定までである。所在の特定が「鍵」に変わるのは、三つ組基準(§8.1)の「予言」項が定量化され、観測可能な形に翻訳されたときであり、それは測度の選定(論文16
判定表)と力学線の将来の仕事である。

\subsection{10. 結論}\label{ux7d50ux8ad6}

当初目指していたのは単なるトイモデルだった。現況は §5
の辞書を確認いただきたい ---
運動学的な主要項は全て埋まり、\textbf{全出典が公開され}、五項では標準理論が公理で置くものを定理として上回り、残額は測度と作用の二項に確定している。著者の結論として述べる:\textbf{標準量子理論やゲージ理論の基本要件には、かなり近い水準で接続できたと考える。}
ここで「基本要件」とは運動学的要件(状態空間・観測量・時間・位相・対称化・保存・測定)を指し、量子理論側はこの全域、ゲージ理論側は骨格(構造要件)まで
--- 作用は未構成のまま --- である。そして接続の不完全な点は §9
で解剖した:それは欠陥の表であると同時に、本体系の経験的内容の所在表であり、将来の判定可能性はそこにのみ住む。この結論は
§1
の読み方の下でのみ主張される:新理論ではなく、\textbf{標準理論の構造が、これほど少ない・これほど異質な仮定からも再出現しうる}ことの実証研究として。

\subsection{11.
主張しないこと}\label{ux4e3bux5f35ux3057ux306aux3044ux3053ux3068}

\begin{enumerate}
\def\labelenumi{(\roman{enumi})}
\tightlist
\item
  標準理論の再導出(方法として選ばない)。(ii)
  物理的同一視(中心テーゼ:同じ結果の別写像)。(iii)
  新理論・統一理論の提案(§1)。(iv) 測度の選定(進行中 ---
  本総説は測度不可知)。(v) 力学。(vi) 偏差の実在(§9.4 ---
  所在の特定まで)。(vii) 本総説は新規の数値検証を含まない ---
  全主張は出典の検証に依拠する(引用地図の機械照合は完了済み)。
\end{enumerate}

\subsection{12. 限界と展望}\label{ux9650ux754cux3068ux5c55ux671b}

\begin{enumerate}
\def\labelenumi{\arabic{enumi}.}
\tightlist
\item
  \textbf{測度}(管轄:論文16)---
  決定構造は有限、判定実験は内部に存在(判定表のとおり進行中:m=5@s31
  ほか)。§9 の「予言」項の定量化はここに係る。
\item
  \textbf{作用}(管轄:将来の力学線)--- 記録列言語への力学の再定式化。
\item
  \textbf{同型定理}(管轄:論文14 系統または独立論文)---
  項目辞書を「運動学的断片上の比較関手の充満忠実性」一枚の定理へ昇格する。材料(Pontryagin
  双対・C⊂H・Z₄↪U(1)・B₄・正準規約)は完備、未証明。
\item
  \textbf{連続極限写像}(管轄:測度決着後の構成課題)--- §9
  横断不一致の「工事中」部分。
\item
  \textbf{(解消済み)出典の公開順} --- 論文13〜16 は 2026-06-12 に Zenodo
  公開され、§5
  の全項が「公」となった。本稿はこれをもって出典確定の確定版である(解消の経過を記録として残す)。
\end{enumerate}

\subsection{13. 追記:起源の総括 ---
そして「単なる複素数ではないか」という批判について}\label{ux8ffdux8a18ux8d77ux6e90ux306eux7dcfux62ec-ux305dux3057ux3066ux5358ux306aux308bux8907ux7d20ux6570ux3067ux306fux306aux3044ux304bux3068ux3044ux3046ux6279ux5224ux306bux3064ux3044ux3066}

\subsubsection{13.1
出発点の経緯(著者の総括)}\label{ux51faux767aux70b9ux306eux7d4cux7defux8457ux8005ux306eux7dcfux62ec}

この一連の論文作成の初期は、λ と ν は共役で λ=1/ν
という強い拘束と、エネルギー様に見える保存量が Σνₙ² = N²
という振動数の二乗和として保存される、というシンプルな仮定から出発した。変位
δ がなければ構造が現れないことは初期から分かっていたので、δ
はほぼ最初から導入していた。検討の途中で、ν と λ が直交している可能性
--- すなわち sin 波と cos 波の関係にあること --- に気づき、さらに δ=±1/2
であることにも気づいた。

この時点で、\textbf{我々が扱っているのは実は、正準化して振幅を1に固定した、複素共役対をなす複素数の極座標表現にすぎない可能性}にも気づいた。しかし本シリーズでは、その方向での再導出はしていない。むしろ、最後まで振幅を持たない(振幅固定1の)ν
と λ として扱い続けたことで、気づけた事象が多くあったと考えている。

\subsubsection{13.2
「単なる複素数をもっともらしく書いただけではないか」への回答}\label{ux5358ux306aux308bux8907ux7d20ux6570ux3092ux3082ux3063ux3068ux3082ux3089ux3057ux304fux66f8ux3044ux305fux3060ux3051ux3067ux306fux306aux3044ux304bux3078ux306eux56deux7b54}

この批判は予想されるので、先回りして四点で答える。

\begin{enumerate}
\def\labelenumi{\arabic{enumi}.}
\item
  \textbf{指摘の核は承認し、開示する。} 振幅1に固定された cos/sin
  対と位相のなす構造が、単位円 U(1)⊂ℂ
  とその共役対の極座標表現に局所的に対応すること ---
  これは事実であり、著者ら自身が検討の途中で気づいた。隠すどころか、ここに明記する。
\item
  \textbf{しかし導出の向きが逆である。} 公理表(§3)を監査されたい ---
  複素数体 ℂ
  もその共役操作も\textbf{入力に存在しない}。複素構造は出力である:C⊂H
  は時間方向の選択から定理として出現し(§5)、共役は反転対合として内部に出現した(§7-2)。「結果が複素数に見える」なら、それは\textbf{複素構造が関係1・定数1・1ビットから再出現した}ことの確認であって、本シリーズの主張の失敗ではなく成功例そのものである。批判が刺さるのは
  ℂ が入力に混入している場合に限られ、公理表は短く、監査可能である。
\item
  \textbf{同型は部分的であり、「単なる複素数」と呼ぶには足りないものと余るものが多すぎる。}
  足りないもの:振幅の自由度(体系には存在しない ---
  固定1)、連続位相(状態データとして現れる位相は連続 U(1) でなく導出
  Z₄)。複素「平面」は状態空間のどこにも現れず、複素線係数は記録の読み出しに導出物として現れるだけである。余るもの:νλ=1
  の格子算術、殻計数(Σνₙ²=N² は周波数空間の格子点の勘定であり、そこから
  1, 9, 137 が出る)、δ=±1/2
  の定数としての地位、パリティ・枝・読み出し階層の構造。極座標の複素数と呼んだ瞬間に、これらは全て表現の影に隠れる。
\item
  \textbf{表現の選択は発見的に非自明だった ---
  そして批判の厳密版は、本シリーズの次の定理である。}
  早期に「複素数の極座標」へ正準化していれば、上記の格子算術と計数構造は「自明な書き換え」として見過ごされた可能性が高い。同値な表現の間でも、何が見えるかは表現に依存する。振幅なしの
  ν/λ 表現を最後まで貫いたことが、§5〜§7
  の発見群に至った経路である。そして「結局は複素数の言い換えではないか」を\textbf{厳密に}述べようとすれば、それは「本体系の運動学的断片と複素ヒルベルト空間の断片との間の構造保存対応を構成し、その充満忠実性を証明せよ」という課題になる
  --- すなわち §12-3
  の\textbf{同型定理そのもの}である。批判を厳密化すると、本シリーズが次に証明しようとしている定理が得られる。その意味でこの批判は、敵ではなく工程表である。
\end{enumerate}

\subsection{参考文献(シリーズ内)}\label{ux53c2ux8003ux6587ux732eux30b7ux30eaux30fcux30baux5185}

\textbf{方針}:本シリーズは外部文献を意図的に引用しない。本稿の主張は構造水準の対応であり、優先権や学説史上の系譜の主張を含まないためである(再構成プログラムへの言及(§1)は一般論であり特定文献への依拠ではない)。引用は全てシリーズ内の公開物と公開リポジトリである。

\textbf{公開済み論文(Zenodo、Concept DOI、全16本)}:

{\def\LTcaptype{none} % do not increment counter
\begin{longtable}[]{@{}lll@{}}
\toprule\noalign{}
\# & タイトル & Concept DOI \\
\midrule\noalign{}
\endhead
\bottomrule\noalign{}
\endlastfoot
1 & 波長空間と周波数空間の双対幾何 & 10.5281/zenodo.20588036 \\
2 & 4D格子セル計数の半径掃引 & 10.5281/zenodo.20588038 \\
3 & 閉4D構造と格子計数(論文1の追補) & 10.5281/zenodo.20589261 \\
4 & 相反双対セル分解と階層真空 & 10.5281/zenodo.20638962 \\
5 & 波の辞書と R 量子化 & 10.5281/zenodo.20640454 \\
6 & 保存則・記録・時間 & 10.5281/zenodo.20640456 \\
7 & 配置統計と関係限定読み出し & 10.5281/zenodo.20640458 \\
8 & 二つの会計と内部膨張・面積則 & 10.5281/zenodo.20640460 \\
9 & 論理波と検閲・安定性 & 10.5281/zenodo.20640462 \\
10 & 創発ゲージ構造 & 10.5281/zenodo.20640464 \\
11 & 次元の必然性 & 10.5281/zenodo.20640466 \\
12 & 双子最小空間モデル & 10.5281/zenodo.20640468 \\
13 & 位置は読まれていなかった & 10.5281/zenodo.20665633 \\
14 & 時間方向と振幅の舞台 & 10.5281/zenodo.20665661 \\
15 & xyztRQ への射影 & 10.5281/zenodo.20665688 \\
16 & 測度問題の厳密な縮約 & 10.5281/zenodo.20665699 \\
\end{longtable}
}

\textbf{補遺とアクセス方法}:補遺(84本)・全検証スクリプト・図の生成コード・二者検算の往復記録・論文13〜16
の査読記録(計8査読)・本稿草稿群は、公開リポジトリの本シリーズフォルダにある。\textbf{正式
URL(コミット時点固定)}:

\begin{quote}
https://github.com/WurabeSeiji/ai-chat-logs-open/tree/547af3c5c005fffc646ecf6ab513a200a3956ed2/\%E6\%B3\%A2\%E9\%95\%B7\%E7\%A9\%BA\%E9\%96\%93\%E3\%81\%A8\%E5\%91\%A8\%E6\%B3\%A2\%E6\%95\%B0\%E7\%A9\%BA\%E9\%96\%93\%E3\%81\%AE\%E5\%8F\%8C\%E5\%AF\%BE\%E5\%B9\%BE\%E4\%BD\%95
\end{quote}

上記はコミット 547af3c5 に固定された時点参照であり、\textbf{URL
自体が日付つきの不変記録}として機能する(§2「記録の所在」の技術的裏づけ)。論文13〜16
公開を含む\textbf{公開作業完了時点の固定参照}(本稿
v1.0・全16論文・索引・スナップショット DOI 記入を含む、コミット
0e96abb2):

\begin{quote}
https://github.com/WurabeSeiji/ai-chat-logs-open/tree/0e96abb2ea995e318b0ccea9d940d230abbaabc3/\%E6\%B3\%A2\%E9\%95\%B7\%E7\%A9\%BA\%E9\%96\%93\%E3\%81\%A8\%E5\%91\%A8\%E6\%B3\%A2\%E6\%95\%B0\%E7\%A9\%BA\%E9\%96\%93\%E3\%81\%AE\%E5\%8F\%8C\%E5\%AF\%BE\%E5\%B9\%BE\%E4\%BD\%95
\end{quote}

以後の最新状態は同リポジトリ main
ブランチの同名フォルダを参照。索引は同フォルダの『論文一覧.md』(全16本の
DOI・読み方3ルート)および『補遺一覧.md』(補遺84本の内容・収載先対応)。補遺1〜41
は論文5〜12 に、42〜48 は論文13 に、49〜61 は論文14 に、62〜67 は論文15
に、68〜84 は論文16
に収載済み(測度の判定実験は進行中)。\textbf{リポジトリのスナップショット
DOI}: 10.5281/zenodo.20666114(Concept;本版 10.5281/zenodo.20666115 ---
2026-06-12、コミット
6fae9b56、論文1〜16・補遺1〜84・検証スクリプト・査読記録の全篇、339MB)。上記
URL と併記して恒久引用に用いる。

\textbf{シリーズ外参照(1件)}:§9.2 の α 恒等式(計数恒等式 α⁻¹=137+(π²/2)α
型、Zenodo 公開済み --- DOI は引用照合の補記に従う)。

\subsection{付録A:主張↔出典↔公開状態↔検証方式}\label{ux4ed8ux9332aux4e3bux5f35ux51faux5178ux516cux958bux72b6ux614bux691cux8a3cux65b9ux5f0f}

§5 の表がこれを兼ねる(全出典公開済み・DOI
は参考文献)。検証方式(全数/標本/乱数)の詳細は各出典の付録に従う。引用地図の機械照合は
Claude Code により完了(2026-06-12)。リポジトリ・スナップショット DOI
は取得・記入済み(10.5281/zenodo.20666114)。

\textbf{謝辞・手続き}:本総説は Claude Code(ローカル機械検証)と
claude.ai(独立検算・査読・本稿起草)の二者検算プロトコルの下で実施・作成された。物理的同一視は行わない。

\end{document}
